Let \(f: [0,1]\rightarrow \mathbb{R}\) continuous. We say that \(f\) crosses the axis at \(x\) if \(f(x)=0\) but \(\exists y,z \in [x-\epsilon,x+\epsilon]: f(y)<0<f(z)\) for any \(\epsilon\).
\begin{itemize}

\item[(a)] Give an example of a function that crosses the axis infinitely often.

\item[(b)] Can a continuous function cross the axis uncountably often?
\end{itemize}
